% ભાગ: ગણો-વિધેયો-સંબંધો
% પ્રકરણ: અંકગણિતીકરણ
% વિભાગ: સંમેય સંખ્યાઓ

\documentclass[../../../include/open-logic-section]{subfiles}

\begin{document}

\olfileid{sfr}{arith}{rat}
\olsection{$\Int$થી $\Rat$ સુધી}

હમણાં જ આપણે સહજ ગણસિદ્ધાંતનો ઉપયોગ કરીને પ્રાકૃતિક સંખ્યાઓમાંથી
પૂર્ણાંકો કેવી રીતે રચવા તે જોયું. હવે તદ્દન સમાન રીતે પૂર્ણાંકોમાંથી
સંમેય સંખ્યાઓ કેવી રીતે રચવી તે જોઈશું. આપણાં પ્રારંભિક અવલોકનો આ છે:
\begin{enumerate}
  \item દરેક સંમેય સંખ્યાને $\nicefrac{i}{j}$ સ્વરૂપે લખી શકાય છે,
  જ્યાં $i$ અને $j$ બંને પૂર્ણાંક છે, પરંતુ $j$ શૂન્યેતર છે.
  \item $\nicefrac{i}{j}$ પદાવલિમાં રહેલી માહિતીને ક્રમયુક્ત જોડ
  $\tuple{ i, j}$માં પણ બરાબર નિરૂપિત કરી શકાય છે.
\end{enumerate}
સહજ અભિગમ સંમેય સંખ્યાઓને $\Int \times (\Int \setminus \{0_\Int\})$
માંથી લેવાયેલી ક્રમયુક્ત જોડો \emph{તરીકે} વિચારવાનો હશે. જોકે પહેલાંની
જેમ આ થોડું અતિસરળ થશે, કારણ કે આપણે $\nicefrac{3}{2} = \nicefrac{6}{4}$
ઇચ્છીએ છીએ, પરંતુ $\tuple{ 3, 2}\neq \tuple{ 6, 4}$. વધુ સામાન્ય રીતે,
આપણને નીચેનું જોઈએ છે:
\[
  \nicefrac{a}{b} = \nicefrac{c}{d} \text{ ત્યારે અને માત્ર ત્યારે જ જ્યારે } a \times d = b \times c
\]
આ મેળવવા માટે આપણે $\Int \times (\Int \setminus \{0_\Int\})$ પર
સામ્ય સંબંધની નીચે મુજબ વ્યાખ્યા કરીએ છીએ:
\[
  \tuple{ a, b }\Ratequiv \tuple{ c, d} \text{ ત્યારે અને માત્ર ત્યારે જ જ્યારે }a \times d = b \times c
\]
આ સામ્ય સંબંધ છે તેની તપાસ કરવી પડશે. આ તપાસ $\Intequiv$ના કિસ્સા
જેવી જ છે, તેથી તેને સ્વાધ્યાય તરીકે રાખીએ છીએ.
\begin{prob}
$\Ratequiv$ એ સામ્ય સંબંધ છે તે બતાવો.
\end{prob}
પરંતુ આ સંબંધ આપણને નીચેની વ્યાખ્યા આપવા દે છે:
\begin{defn}
સંમેય સંખ્યાઓ એ $\Ratequiv$ હેઠળ પૂર્ણાંકોની એવી જોડોના સામ્ય વર્ગો છે
જેનો બીજો ઘટક શૂન્યેતર હોય. એટલે કે, $\Rat =
\equivclass{(\Int \times (\Int\setminus \{0_\Int\}))}{\Ratequiv}$.
\end{defn}

પૂર્ણાંકોની જેમ અહીં પણ કેટલીક મૂળભૂત ક્રિયાઓ વ્યાખ્યાયિત કરવી છે.
$\equivrep{i,j}{\Ratequiv}$ એ $\Ratequiv$ હેઠળનો એવો સામ્ય વર્ગ હોય
જેમાં $\tuple{i, j}$ !!a{element} છે, ત્યારે આપણે કહીએ છીએ:
\begin{align*}
  \equivrep{a, b}{\Ratequiv} + \equivrep{c, d}{\Ratequiv} &= \equivrep{ad + bc,  bd}{\Ratequiv}\\
  \equivrep{a, b}{\Ratequiv} \times \equivrep{c, d}{\Ratequiv} &= \equivrep{a  c, b d}{\Ratequiv}.
\intertext{આ સંમેય સંખ્યાઓ પર $r \leq s$ની વ્યાખ્યા કરવા માટે આપણે એ
હકીકત વાપરીએ છીએ કે $r \le s$ ત્યારે અને માત્ર ત્યારે જ જ્યારે
$s - r$ ઋણ ન હોય; એટલે કે $s - r$ને $\nicefrac{i}{j}$ સ્વરૂપે લખી
શકાય, જ્યાં $i$ અનૃણ અને $j$ ધન હોય:}
  \equivrep{a, b}{\Ratequiv} \leq \equivrep{c, d}{\Ratequiv} &\text{ ત્યારે અને માત્ર ત્યારે જ જ્યારે }
  \equivrep{c, d}{\Ratequiv} - \equivrep{a, b}{\Ratequiv} =
  \equivrep{i_\Int, j_\Int}{\Ratequiv}
\end{align*}
કોઈ $i \in \Nat$ અને $0 \neq j \in \Nat$ માટે.
\sourcecorrection{OLARI-001}{મૂળની \texttt{rationals.tex}, પંક્તિઓ
53--58માં અસમાનતા અને દર્શાવેલું સૂત્ર બંને \(s-r\) અનૃણ હોવું જોઈએ
એવું કહે છે, પરંતુ વચ્ચેના વાક્યમાં ભૂલથી \(r-s\) લખાયું છે. અહીં
સૂત્ર અને અભિપ્રાય સાથે સુસંગત \(s-r\) વાપર્યું છે.}

પછી આ વ્યાખ્યાઓનો વ્યવહાર જેવો \emph{હોવો જોઈએ} તેવો જ છે તેની
તપાસ કરવી પડે છે; આ તપાસ \olref[check]{sec}માં રાખી છે. પરંતુ તેમનો
વ્યવહાર ખરેખર એવો જ છે!{} છેલ્લે પૂર્ણાંકોને સંમેય સંખ્યાઓ
\emph{તરીકે} ગણવાની કોઈ રીત જોઈએ; તેથી દરેક $i \in \Int$ માટે આપણે
$i_\Rat = \equivrep{i, 1_\Int}{\Ratequiv}$ એમ ઠરાવીએ છીએ. આ બધું
યોગ્ય રીતે વર્તે છે તેની તપાસ પણ \olref[check]{sec}માં કરીએ છીએ.

\begin{prob}
કોઈ પણ $i, j \in \Int$ માટે $(i + j)_\Rat = i_\Rat+ j_\Rat$,
$(i \times j)_\Rat = i_\Rat \times j_\Rat$ અને
$i \leq j \liff i_\Rat \leq j_\Rat$ છે તે બતાવો.
\end{prob}

\end{document}
